%! Polynomial Functions and Complex Zeros — Unit 2, Lesson 2.4 — Student Packet Cover Sheet
\documentclass[10pt]{article}
\usepackage{precalculus-article}
\usepackage{precalculus-boxes}
\usepackage{ltablex}
\keepXColumns

\begin{document}

% Full-bleed plum banner
\begin{tikzpicture}[remember picture, overlay]
  \fill[plum] (current page.north west)
    rectangle ([yshift=-0.9in]current page.north east);
\end{tikzpicture}
\vspace{-0.8in}

\begin{center}
  {\color{white}\textbf{\LARGE Precalculus}}\\[4pt]
  {\color{lilacmid}\large\textbf{Unit 2: Polynomial Functions}}\\[6pt]
  {\color{white}\textbf{\large Lesson 2.4 \quad Polynomial Functions and Complex Zeros}}
\end{center}

\vspace{0.22in}
\namedateperiod
\vspace{0.18in}

\begin{learningtargetbox}
\small
\begin{enumerate}[leftmargin=1.5em, itemsep=5pt]
  \item I can explain why a polynomial can have zeros that never appear as
        $x$-intercepts, and I can evaluate a polynomial at a complex input to
        check whether that input is a zero.
  \item I can use the Fundamental Theorem of Algebra to say that a polynomial
        of degree $n$ has exactly $n$ complex zeros counting multiplicities,
        and I can split that total into the real zeros a graph shows and the
        non-real zeros it hides.
  \item I can use conjugate pairs to name a missing non-real zero, explain why
        an odd-degree polynomial must have a real zero, and build a polynomial
        from a complete list of its zeros.
  \item I can decide whether a polynomial function is even, odd, or neither,
        both by substituting $-x$ and by reading the symmetry of a graph.
\end{enumerate}
\end{learningtargetbox}

\vspace{0.14in}

\begin{tocbox}
\small
\renewcommand{\arraystretch}{1.5}
\begin{tabularx}{\linewidth}{@{} c l X r @{}}
\rowcolor{lilac}
\textbf{\#} & \textbf{Component} & \textbf{Description} & \textbf{Score} \\
\midrule
1 & Warm-Up        & Spiral review of prerequisite skills                  & \blank{1.2cm} \\
2 & Guided Notes   & Complex zeros, conjugate pairs, even and odd functions & \blank{1.2cm} \\
3 & Group Activity & The Missing Zeros Bureau: accounting for every zero    & \blank{1.2cm} \\
4 & Exit Ticket    & Independent check on counting and classifying zeros    & \blank{1.2cm} \\
5 & Homework       & Practice and extension                                 & \blank{1.2cm} \\
\midrule
  & \multicolumn{2}{r}{\textbf{Total}} & \blank{1.2cm} \\
\end{tabularx}
\end{tocbox}

\vspace{0.14in}

\begin{remindbox}
\small
Last lesson every zero was something you could point at on the graph. Today
the count stops depending on what you can see: a degree-$n$ polynomial has
\textbf{exactly} $n$ zeros, and the ones the graph refuses to show up as
$x$-intercepts are the non-real ones --- always arriving two at a time, always
as a conjugate pair.
\end{remindbox}

\end{document}
